\Qhead{2.7}{Метод Кирхгофа в теории дифракции}{Метод Кирхгофа в теории дифракции. Функции Грина. Условия излучения. Дифракция в зоне Френеля и Фраунгофера. Характеристики поля в фокусе линзы.}
\Subfile{2.7.1}{sections/02_waves/q07/s01}
\Subfile{2.7.2}{sections/02_waves/q07/s02}

\recap{Функция Грина превращает волновую задачу в граничный интеграл. Кирхгоф приближённо задаёт поле на отверстии; Френель сохраняет квадратичную фазу, Фраунгофер даёт преобразование Фурье. Линза реализует преобразование Фурье в фокальной плоскости.}
\board{Исходящая $G$; интеграл Френеля; $N_F$; щель или пятно Эйри.}
\oral{Почему условие излучения необходимо дополнительно к уравнению Гельмгольца?}
\bait{Можно оставить роль линзы: её квадратичная фаза позволяет получить дальнюю дифракционную картину на конечном расстоянии.}
\sources{Warde~\cite{wardediffraction}, лабораторная работа~3; Fang~\cite{fangcoherence}, раздел дифракции и фурье-оптики. Нормировка граничного и параксиального интегралов дополнительно проверена прямым выводом; в пособии Белокопытова не заявляется несуществующий раздел дифракции.}
